\section[两个拓扑场论的故事]{两个拓扑场论的故事（1 月 21 日）}
\srcpage{5}

上次我们谈了模形式（或自守形式）的一个视角：取你最喜欢的约化代数群或矩阵群，例如 $\mathrm{GL}_n$ 或 $\mathrm{PSL}_2$\footnote{我们并不是把这些群具体地看作某个特定域上的群，例如 $\mathrm{GL}_n(\mathbb{R})$，而是把它们看作一台机器：给一个域 $k$，就给出一个群 $\mathrm{GL}_n(k)$。今天要让 $F$ 变动，这个技术细节很重要。}，再令 $F/\mathbb{Q}$ 为数域。愿意的话你可以取 $F=\mathbb{Q}$。设 $\mathcal{O}_F$ 为 $F$ 的整数环；若 $F=\mathbb{Q}$，则 $\mathcal{O}_F=\mathbb{Z}$。

今天我们要讨论改变 $F$ 时会发生什么，以及这如何影响主丛的模空间（TODO：这里漏听了）。我们这样得到局部对称空间：取群的实点，再对「算术」格与极大紧子群取双商。例如可以取 $\mathrm{PSL}_2(\mathbb{Z})\backslash\mathrm{PSL}_2(\mathbb{R})/\mathrm{SO}_2$。这里取 $S=\mathbb{Q}$。

现在考虑更一般的 $F$。我们有 $\mathrm{PSL}_2(\mathcal{O}_F)$，这没问题，但该用什么替换 $\mathrm{PSL}_2(\mathbb{R})$ 及其极大紧子群？我们改为考虑 $F\otimes_{\mathbb{Q}}\mathbb{R}$，这是一个形如
\[
F\otimes_{\mathbb{Q}}\mathbb{R}\cong\mathbb{R}^{r_1}\times\mathbb{C}^{r_2}
\]
的环，其中 $F$ 有 $r_1$ 个映入 $\mathbb{R}$ 的嵌入，以及 $r_2$ 对共轭的映入 $\mathbb{C}$ 的嵌入。于是我们把 $\mathrm{PSL}_2(\mathbb{R})$ 换成 $\mathrm{PSL}_2(F\otimes_{\mathbb{Q}}\mathbb{R})$。

\begin{example}
设 $F=\mathbb{Q}(\sqrt{d})$，其中 $d$ 无平方因子。
\begin{itemize}
  \item 若 $d>0$，这是实二次扩张，$r_1=2$、$r_2=0$。
  \item 若 $d<0$，这是虚二次扩张，$r_1=0$、$r_2=1$。
\end{itemize}
\end{example}

然后可以照常取 $\mathrm{PSL}_2(F\otimes_{\mathbb{Q}}\mathbb{R})$ 的极大紧子群 $K$，得到一个局部对称空间。若 $F$ 是实二次域，就经由 $\mathrm{PSL}_2(\mathbb{R})\times\mathrm{PSL}_2(\mathbb{R})$ 作用在 $\mathbb{H}\times\mathbb{H}$ 上，得到 Hilbert 模形式。在虚二次的情形，我们得到 $\mathrm{PSL}_2(\mathbb{C})/\mathrm{SO}_3\cong\mathbb{H}^3$，即双曲 3-空间，与双曲几何有联系。

\srcpage{6}

通常这些双商不是代数簇，这个例子就说明了这一点。

对其他 $F$，我们会得到 $\mathrm{PSL}_2(\mathbb{R})$ 与 $\mathrm{PSL}_2(\mathbb{C})$ 的乘积；最有意思的是算术格 $\mathrm{PSL}_2(\mathcal{O}_F)$。

一旦有了这些算术局部对称空间 $M$，我们想从它们造出向量空间，包括 $L^2(M)$、$H^*(M)$ 及其扭版本。对几何 Langlands 纲领而言重要的是：这些向量空间带有一个巨大的交换代数的作用，而它是所有素元上的多项式环的张量积，每个多项式环有 $\operatorname{rank}(G)$ 个生成元。

也可以允许分歧，得到推广 $M_{G,F,N}$，其中 $N\in\mathcal{O}_F$，并把 $\mathrm{PSL}_2(\mathcal{O}_F)$ 换成同余子群 $\Gamma_N$；后者是对矩阵模 $N$ 施加某些条件。你可以施加的条件有几种（例如模 $N$ 等于恒等，或模 $N$ 为上三角）。算术局部对称空间是 $\Gamma_N\backslash G_{\mathbb{R}}/K$，而那个大交换代数现在「只是」对不整除 $N$ 的素数 $p$ 取张量积。

这类想法——把一个向量空间联系到数域，或者联系到数域连同一些素元——在算术拓扑的类比下，让人想起 4 维拓扑场论中的状态空间。如上次所讨论，这是 Weil 的 Rosetta 石的一个精细化：$\Spec\mathbb{Z}$ 感觉像一条曲线，其上的点 $\Spec\mathbb{F}_p$ 对应于素数 $p$，而 $\Spec\mathbb{Z}_p$ 是这个点周围的一个小圆盘。圆盘里面是去心圆盘 $\Spec\mathbb{Q}_p$。这类似于有限域 $\mathbb{F}_q$ 上一条光滑约化的代数曲线，它在局部看起来像 $\Spec\mathbb{F}_q[t]$。这里的点是 $\Spec\mathbb{F}_q$，点周围是圆盘 $\Spec\mathbb{F}_q[[t]]$，以及去心圆盘 $\Spec\mathbb{F}_q((t))$。

现在我们来看 $\Spec\mathbb{Z}$ 的 étale 拓扑；这是一种花哨的说法，意思是我们关心 Galois 群的上同调。从这个视角看，Rosetta 石还不够丰富：$\Spec\mathbb{F}_p\hookrightarrow\Spec\mathcal{O}_F$ 是一个点，但从 étale 拓扑的意义上说并不是：étale 意象告诉我们，这个「点」有一大堆有意思的覆叠空间，例如 $\Spec(\mathbb{F}_{p^n})\to\Spec(\mathbb{F}_p)$。它的 étale 基本群是 $\widehat{\pi}_1^{\text{ét}}(\Spec\mathbb{F}_p)=\Gal(\overline{\mathbb{F}}_p)\cong\widehat{\mathbb{Z}}$。

你可以认为这种投射有限性其实并不存在：我们只能看见有限扩张，或者换个说法，只能看见有限覆叠空间；在这个层面上，没有办法区分 $\mathbb{Z}$ 与 $\widehat{\mathbb{Z}}$。这是 étale 基本群的一个常见特点。

所以要点是：如果你想做 étale 的事情，点 $\Spec\mathbb{F}_p$ 的行为很像一个圆。（TODO：画 $\Spec\mathbb{Z}$，每个素元处画一个圆。）因此 $\Spec\mathbb{Z}$ 及其同类 $\Spec\mathcal{O}_F$ 感觉上更像 3-流形，而不是 Riemann 曲面。还有别的办法让这个模糊的类比不那么模糊：例如存在一个 Poincaré 对偶的版本，维数也正确。

绕过这些圆的单值是 Frobenius，但不同素数处的不同 Frobenius 彼此不说话。因为曲线 $C/\mathbb{F}_q$ 映到 $\Spec(\mathbb{F}_q)$，而后者有点像圆，我们就把这些 3-流形想成纤维化在圆上的。

在这个视角下，$\Spec(\mathbb{F}_q)$ 是什么？从 étale 拓扑看，它其实是一个点；但若你想要覆叠空间与 Galois 表示之间的好字典，它应当是圆的一个覆叠空间，而且事实上是万有覆叠 $\mathbb{R}$。这没问题：就拓扑与上同调的目的而言，$\mathbb{R}$ 是点的很好替身。现在把 $C$ 基变换为 $\overline{C}:=C_{\overline{\mathbb{F}}_q}:=C\times_{\Spec\mathbb{F}_q}\Spec\overline{\mathbb{F}}_q$；此时我们得到的东西感觉像 Riemann 曲面丛，即代数闭域上的曲线，纤维化在实直线 $\Spec\overline{\mathbb{F}}_q$ 上。

如上次所讨论，绕过圆的单值是 Frobenius，所以我们可以把这些曲面丛的类比物想成 Frobenius 的映射环面。

再讨论一条支持这个算术拓扑字典的证据：$\Spec\mathbb{Z}$ 会怎样？设 $p$ 为素数，于是得到 $\Spec\mathbb{Z}$ 中的一个「圆」$\Spec\mathbb{F}_p$。邻域 $\Spec\mathbb{Q}_p$ 现在表现得像这个圆在 $\Spec\mathbb{Z}$ 内部的管状邻域。更一般地，可以取 $\Spec\mathcal{O}_F$、素元 $p\in F$ 以及一个位 $v$，把 $\mathcal{O}_F$ 在 $v$ 处完备化，得到局部域 $F_v$；于是我们可能预期 $\Spec(F_v)$ 是圆 $\Spec\mathbb{F}/p$ 的管状邻域——不过（TODO）这个位得知道一些关于 $p$ 的信息。

若 $F_v$ 是非阿基米德局部域，例如 $\mathbb{Q}_p$ 或 $\mathbb{F}_p((t))$，则 $\Gal(\overline{F}_v/F_v)$ 满射到 $\mathbb{Z}_\ell\rtimes\widehat{\mathbb{Z}}$ 上。给对群论感兴趣的人一个有趣的事实：这个半直积是 Baumslag–Solitar 群
\[
\mathrm{BS}(1,p):=\langle \sigma,u\mid \sigma u\sigma=u^p\rangle
\]
的一个例子。这里 $\sigma$ 是 Frobenius，$u$ 是 $\mathbb{Z}_\ell$ 的生成元。它在 $p=1$ 时的 $\mathbb{Z}\times\mathbb{Z}=\pi_1(T^2)$ 与 $p=-1$ 时的 Klein 瓶的基本群之间做插值。所以这确实有点像 3-流形中一个纽结的环面邻域。

\srcpage{7}

于是素数看起来像圆，局部域看起来像 2-流形——还不是随便什么 2-流形，而是纤维化在圆上的 2-流形。

一般来说，Galois 群 $\Gal(\overline{F}_v/F_v)$ 可以由三部分拼成：剩余域的 Galois 群 $\widehat{\mathbb{Z}}\cdot\sigma$；驯顺部分（tame part），它是 $\ell\ne p$ 的那些 $\mathbb{Z}_\ell$ 的乘积（这里 $p$ 是剩余域的特征）；以及野部分（wild part），它是一个 $p$-群。

这个类比很好也很重要，它把算术 Langlands 与几何 Langlands 绑在一起；但我们会把大部分时间花在它最具体的地方。让我们把这个类比总结一下。

\begin{itemize}
  \item 以下对象被看作三维的：数域 $\Spec\mathcal{O}_F$ 与曲线 $C/\mathbb{F}_q$ 的函数域，以及 Riemann 曲面微分同胚 $\varphi:\Sigma\to\Sigma$ 的映射环面。前两者与整体 Langlands 纲领有关，我们称之为「整体算术设定」。
  \item 以下是一些二维对象：局部域 $F_v/\mathbb{Q}_p$ 及其谱，它们像去心圆盘；还有 $\mathbb{F}_q((t))$，它也是某种去心圆盘。这是「局部算术设定」中局部 Langlands 纲领的舞台。还有另外两种二维对象：正特征代数闭域上的曲线 $C/\overline{\mathbb{F}}_q$，或闭 Riemann 曲面 $\Sigma$。后两种对象构成「整体几何设定」。
  \item 一维对象：$\Spec\mathbb{F}_q$ 与 $\Spec\mathbb{C}((t))$ 都是圆的类比物。后者是去心圆盘，所以不恰好是一维的，但已经够接近、够用了；它是「局部几何设定」。
  \item 最后是零维对象：$\Spec\mathbb{F}_q$ 与 $\Spec\mathbb{C}$。
\end{itemize}

本课程的一个主要主题，是把一个 4 维拓扑场论应用到这些对象上。如果你抱怨这里根本没有 4-流形，那是个好问题；但我们只会考虑少数几个 4-流形，例如 3-流形与圆的乘积，或者更一般的映射环面。

Langlands 纲领是 4 维「算术拓扑场论」之间的一个等价。算术 TFT 并不是一个完全良定义的对象，但这样一份定义所需的大部分数据我们都有。取一个群 $G$，得到对偶群 $G^\vee$；于是这两个算术 TFT 称为 $A_G$ 与 $B_{G^\vee}$：$A_G$ 称为自守侧或磁侧，$B_{G^\vee}$ 称为谱侧或电侧。某种意义上，这是四维版本的镜像对称（通常的镜像对称是二维量子场论的故事）。这两个 TFT $A_G$ 与 $B_{G^\vee}$ 都是全扩张的，也就是说，它们把高阶范畴对象指派给低维对象。具体地说，在上述字典里，我们可以把东西指派给二维、一维、也许还有零维对象：二维流形得到范畴，一维流形得到 2-范畴，如果你非常有雄心，零维流形得到 3-范畴。

$A_G$ 是一台机器：取 3-流形 $M$，附上一个向量空间 $A_G(M)$；我们将会看到，它由算术局部对称空间上的函数造出。例如 $\Spec\mathcal{O}_F$ 得到 $M_{G,F}$ 上某种函数空间。这是很大的结构，好处之一是它能解释模形式的一些古怪性质。我们也会花些时间在 $B$ 侧，那一侧更容易描述。

维数、范畴层数与难度之间有一个有趣的权衡：要弄清这些算术 TFT 给 4-流形指派什么，非常困难——要处理无穷大与棘手的重整化，还有超出本课程范围的分析。向量空间更美好、更容易造，但 3-流形很难。维数 2 是最佳平衡点：范畴还不算太糟，2-流形也相当好处理。等走到 1-流形，就得把范畴论用得更狠；而要理解 0 维发生的事情，几乎完全是未知的。我们不会在本课程中解决这个未解问题。

TFT 还有额外的结构：你可以取一个带额外结构的流形，这种结构称为缺陷（defects），并给它也指派代数数据。这些附加装置在算术侧与拓扑侧对得非常齐：例如，数论学家会告诉你，在定义算术局部对称空间时允许分歧/同余子群是多么重要。在算术拓扑字典下，这对应于研究你的 TFT 给一个带嵌入链环（链环带标号）的 3-流形指派什么。用物理的语言，这个链环是余维 2 缺陷。链环的额外数据给出用同余子群造出的修改后的向量空间。

作用在 $M$ 上函数空间的那个大交换代数，含有称为 Hecke 算子的元素；这些对应于余维 3 缺陷，即线算子。从物理上说，你可以把它们想成「产生磁单极子」。而自守形式的周期对应于边界条件——这是一种余维 1 的现象。

\srcpage{8}

这与教授近期的工作以及进行中的工作有关！

最后是 Langlands 函子性，它也能放进这幅图景：不只是边界，还有称为畴壁（domain walls）的余维 1 现象，你可以把它想成流形上两个区域之间的界面，两个区域上承载着两种不同的理论。

所以总结一下：自守形式理论中所有这些附加装置，都属于这个量子场论的故事。

\paragraph{$B$ 侧}

这里我们把自守形式的理论送过一块棱镜，把它分解成与 Galois 表示有关的「颜色」。在数论上，这一侧非常非常难，因为数域的 Galois 群很复杂；而 $A$ 侧常常被用来获取关于 $B$ 侧的信息。在几何上，Riemann 曲面的基本群容易得多，所以 $B$ 侧被用来了解 $A$ 侧。

但至少 $B$ 侧更容易陈述：我们研究 Galois 表示空间的代数几何；相比之下，$A$ 侧与算术局部对称空间 $M$ 的拓扑有关。在几何情形，我们可以固定一个流形 $M$，考虑表示 $\pi_1(M)\to\mathrm{GL}_n(\mathbb{C})$ 构成的空间 $\operatorname{Loc}_n(M)$。这些对象很好，称为特征簇（character varieties），你可以研究其上的函数代数。你也不必把自己限制在 $\mathrm{GL}_n(\mathbb{C})$：我们特别关心 $\operatorname{Loc}_{G^\vee}(M):=\{\pi_1(M)\to G^\vee\}$。$B_{G^\vee}$ 指派给 3-流形的向量空间，就是 $\operatorname{Loc}_{G^\vee}(M)$ 上的函数空间；事实上，把这个理论定义成扩张 TFT，比 $A$ 侧容易得多。回到算术侧，$\pi_1(M)$ 类比于 Galois 群，所以在算术设定中，我们看的是 Galois 表示的簇。

这两个（算术）TFT 的猜想等价性是一件了不起的事：$A$ 侧的巨大结构，等价于更容易定义的 $B$ 侧，所有结构都来回传递。但「猜想」这个词在这里分量很重：在算术与几何两种设定中，都还有很多事情要做，甚至要先定义，才能让这些类比精确。在几何设定中，已知的更多，但仍有许多进行中的工作，包括教授、Arinkin、Gaitsgory、Raskin 以及许多其他人的工作。数域的故事是 Lafforgue 以及许多其他人的工作，但其中被证明的并不多。

这本字典不只是好看：你可以用几何设定中已完成的工作，来了解在算术设定中应当朝什么方向努力，等等。
